% source: https://tex.stackexchange.com/a/763272 (question 724347, score 2, block 1)
\documentclass{standalone}
\usepackage[svgnames]{xcolor}
\usepackage[3d]{luadraw}
\usepackage{fourier-otf}
\begin{document}
\begin{luadraw}{name=sphere}
    local ld = luadraw
    local pt3d = ld.pt3d
    local M = pt3d.M
    local Origin,vecI,vecJ, vecK = pt3d.Origin, pt3d.vecI, pt3d.vecJ, pt3d.vecK
    local sin, cos = math.sin, math.cos
    
    local interBcircles = function(C1, C2, radius)
        -- C1 = {A1, A2}, C2 = {A3, A4} --fours points on the sphere
        local P1 = ld.plane(C1[1],C1[2],M(0,0,0))
        local P2 = ld.plane(C2[1],C2[2],M(0,0,0))
        local D = ld.interPP(P1,P2)
        local u = pt3d.normalize(D[2]) -- director vector of intersection
        return {radius*u, -radius*u}
    end

    require "luadraw_spherical"
    local R = 4
    local g = ld.graph3d:new{window3d = {-6, 6, -6, 6, -6, 6},viewdir = {-40, 70}}
    ld.Hiddenlines = true;ld.Hiddenlinestyle = "dashed"

    g:Define_sphere{
        radius = R,
        color = "white",
        mode = ld.mBorder,
        edgecolor = "black",
        insidelabelcolor = "black!70"
    }
    -- a big circle is defined by two points on the sphere (not aligned with the center).
    local C1, C2, C3 = {R*vecI, R*vecJ}, {R*vecJ, M(-R * cos(35 * ld.deg), 0, R * sin(35 * ld.deg))}, {R*vecI, R*vecK}
    local I1, I2 = table.unpack( interBcircles(C1,C2,R) ) -- intersection between C1 and C2
    local I3, I4 = table.unpack( interBcircles(C1,C3,R) )
    local I5, I6 = table.unpack( interBcircles(C2,C3,R) )

    g:DSbigcircle(C1, {color = "black"})
    g:DSbigcircle(C2, {color = "red"})
    g:DSbigcircle(C3, {color = "green"})
    g:DSstars({I1,I2,I3,I4,I5,I6})
    
    g:Dspherical()
    g:Show()
\end{luadraw}
\end{document}
